\section{The three cut-cocycle conditions over the enlarged coefficient field}

% BEGIN C80L1
\paragraph{C80L1. A presentation of each original cut cocycle after localization.}
Fix a member $D$, a localization $i$, a place $w\mid\ell$, and the
coefficient enlargement $C$ of the native local problem. Write $G$ for
the local absolute Galois group, $L$ for the enlarged coefficient
field, $O$ for its integer ring and $\lambda$ for its chosen
uniformizer. All scalar extensions below use the given localization
map on the original coefficient ring. Write $N_x$ for the localized
integral matrix representation $R_x\otimes\chi_1^{-1}$; for the
cut target write $M$ for the localized $\rho\otimes\chi_1^{-1}$
or $\chi_c\chi_1^{-1}$, in ranks $k=r$ or $1$, respectively.

Let $z$ be any member of the original set $\mathrm{cutCoc}$ for
this target on the full set of points. Write $f$ for its image
in $(L/O)^k$. Then there are a finite list of points $x_j$, a vector
$E\in\prod_j L^{n_V}$, a finitely generated stable $O$-submodule
$Y$ of this product, and an equivariant $O$-linear map
$\phi:Y\to (L/O)^k$ such that, with $N=\bigoplus_j N_{x_j}$,
\[
 d_g=N(g)E-E\in Y,\qquad f(g)=\phi(d_g).
\]
Indeed these are the presentations in the definition of the cut
before localization. Map the entries of $E$ and the original
module to $L$, and take the $O$-span of the image of the module.
A basis of the original finite torsion-free module over its DVR
is independent over its fraction field by clearing denominators.
It remains independent over $L$ by a nonzero minor of its coordinate
matrix, and it generates the new span. Thus the semilinear map obtained
by applying localization and the quotient coefficient inclusion to
the original $\phi$ extends to an $O$-linear map on this span.
On its generators the localized matrices carry the original action
to the new one, so stability and equivariance extend to the entire
span. The specified quotient inclusion sends the class of a vector
to the class of the included vector. This proves the displayed
formula for the original $z$ itself.

There is an $a\ge0$ for which all of $\phi(Y)$ is killed by
$\lambda^a$: take the maximum of the denominators of the finitely
many images of the basis vectors, with zero allowed when the image
is zero. Also $f$ is locally constant. Indeed $Y$ has its ordinary
lattice topology in its $L$-span: the coordinates of its independent
basis are continuous there by a left inverse over $L$. Therefore
$\lambda^a Y$ is open in $Y$ and $\phi$ is locally constant.
Continuity of the chosen matrices in a good family and of the
character twist gives continuity of $g\mapsto d_g$. Finally
$d_{gh}=d_g+N(g)d_h$, so $f(gh)=f(g)+M(g)f(h)$.
These conclusions hold for each generator with its own exponent $a$.
% END C80L1

% BEGIN C80L2
\paragraph{C80L2. Lifting to the exact finite coefficient module and canceling the character.}
Put $X_a=(L/O)[\lambda^a]$. The map
$s\mapsto s/\lambda^a\bmod O$ identifies
$O/\lambda^a O$ with $X_a$. It is finite because the residue field
is finite, and it is cyclic with specified generator
$x_0=\lambda^{-a}\bmod O$. This covers $a=0$, when the module is zero.
In the local data of $\mathrm{SetC}$ let $\theta$ be the shifting
character, put $\kappa_c=\chi_c\chi_1^{-1}$ after localization,
and $\alpha_c=\kappa_c\theta$, so $\alpha_1=\theta$.
Use the standard lattice $T=O^r$ with action $P=\theta M$ in the
rank case, and the standard line $T=O$ with action
$P=\alpha_c$ in the character case. The natural map
\[
 \iota:T\otimes_O X_a\longrightarrow (L/O)^k,\qquad
 t\otimes(s\bmod O)\longmapsto st\bmod O^k
\]
is injective with image exactly $(L/O)^k[\lambda^a]$, by the
standard coordinate basis. On the tensor space give $g$ the action
$\theta(g)^{-1}(P(g)\otimes 1)$. The map $\iota$ then intertwines
it with $M(g)$ in both cases. In particular this is exactly the
native action used to test a finite-coefficient cocycle.

Choose $a$ as in C80L1. There is a unique function $z_a$ with
$\iota z_a=f$. It is locally constant and satisfies the cocycle
law by injectivity. Moreover the map $\phi$ of C80L1 has image in
the image of $\iota$; viewing it there gives an equivariant map
$\phi_a:Y\to T\otimes X_a$ with $z_a(g)=\phi_a(d_g)$.
More generally, if an exponent was chosen to kill just the values
of $f$, replace $Y$ by $\phi^{-1}(\iota(T\otimes X_a))$. This
is stable by equivariance, finitely generated by the Noetherian
property, and still contains every $d_g$. The construction thus
applies to any exponent and any lift as in the original two
admissibility theorems.
% END C80L2

% BEGIN C80L3
\paragraph{C80L3. An explicit module surjecting onto the integral extension.}
Assume the native splitting at $\pi$, and use the notation of
C80L1--C80L2. The local representation $N_\pi$ over $L$ contains
a nonzero fixed vector $e$ on its first-character line. There is
an injective $L$-linear map $j_T:T\otimes L\to L^{n_V}$
intertwining $P$ with $\theta N_\pi$: for the rank block take
the inverse of the specified splitting on the rank summand, and for
the $c$th character take the inverse on the $c$th coordinate line.
These assertions follow by restricting the original splitting to
the decomposition group and mapping its entries and its inverse
to the enlarged field. The first line has ratio character $1$ and the
$c$th line has ratio $\kappa_c$ before the common shift.
When $c=1$ the two vectors may lie on the same line; use different
copies of $N_\pi$ in what follows.

Let $b\ge0$ be any integer killing $X_a$ by $\lambda^b$.
The finite module whose admissibility is required is
\[
 H_b=(T\otimes X_a)\oplus O/\lambda^bO,qquad
 g(m,s)=\big((P(g)\otimes 1)m+s\theta(g)z_a(g),\ \theta(g)s\big).
\]
The choice of lift of $s$ does not affect the correction term.
The cocycle law of C80L2 gives the action law in this formula.
Take $A=Y\oplus T\oplus O$ and embed it in the direct sum of the
$\theta N_{x_j}$ and two further copies of $\theta N_\pi$ via
\[
 J(y,t,s)=\big(y+sE,\ j_T(t),\ se\big).
\]
This is injective since $e\ne0$ and $j_T$ is injective. Its image
is stable, with the action in these coordinates given by
\[
 g(y,t,s)=
 \big(\theta(g)N(g)y+s\theta(g)d_g,\ P(g)t,\ \theta(g)s\big).
\]
Indeed the first component of $J$ after this substitution is
$\theta(g)N(g)(y+sE)$, and the other two follow from the two
intertwining identities. The integral map
\[
 q:A\longrightarrow H_b,\qquad
 (y,t,s)\longmapsto(\phi_a(y)+t\otimes x_0,\;s\bmod\lambda^b)
\]
is surjective, since $x_0$ generates $X_a$. It is equivariant:
in the first coordinate $\phi_a(\theta(g)N(g)y)$ is
$(P(g)\otimes 1)\phi_a(y)$ and
$\phi_a(s\theta(g)d_g)=s\theta(g)z_a(g)$.
Thus the kernel of $q$ is stable and $J(A)/J(\ker q)$ is the
required extension with precisely its original twist and its
second-coordinate modulus $\lambda^b$.
% END C80L3

% BEGIN C80L4
\paragraph{C80L4. Realizing the extension with two full stable lattices.}
By $\mathrm{SetC}$ each of the shifted representations used in
C80L3 is continuous and crystalline, has jump function
$\widetilde j$ prescribed by the position table, and has characteristic
polynomial congruent modulo $\lambda$ to the prescribed product
of the rank-block polynomial and the $t$ character factors.
These are the defining conditions for each summand of the native
$\mathrm{AdmQ}$ witness. Let $V$ be their direct sum, with the
standard full stable product lattice $\Lambda$. Set $A'=J(A)$ and
$B'=J(\ker q)$. Both are stable finitely generated $O$-submodules
of $V$; $\lambda^b A'\subset B'$. They need not span $V$, so
we now construct full lattices yielding exactly the same quotient.

Put $c=\lambda^b\ne0$. In the Noetherian module $\Lambda$ the
ascending chain
$\{v\in\Lambda:c^jv\in A'\}$ stabilizes, say at $j=m$.
With $\mu=c^{m+1}$ this gives
$\mu\Lambda\cap A'\subset cA'\subset B'$:
if $c^{m+1}v\in A'$ then $c^m v\in A'$.
The modules
\[
     \Lambda_1=A'+\mu\Lambda,\qquad
     \Lambda_2=B'+\mu\Lambda
\]
are stable, finitely generated and torsion-free over the DVR,
hence free. They span $V$ because they contain $\mu\Lambda$,
and $\Lambda_2\subseteq\Lambda_1$. The map from $A'$ onto
$\Lambda_1/\Lambda_2$ has kernel
$B'+(A'\cap\mu\Lambda)=B'$. Its induced isomorphism with $H_b$
is equivariant, with the original action on representatives from
C80L3. This proves the actual $\mathrm{AdmQ}$ at $b$, also when
$c=1$. Since every nonnegative $b$ annihilating $X_a$ was allowed,
the cutoff $b_0=0$ proves the original
$\mathrm{IsAdmCoc}$ for both the rank and character extensions,
with all the shifted jumps and congruences on the same finite list
of genuine crystalline representations.
% END C80L4

% BEGIN C80L5
\paragraph{C80L5. The finite part on the rank block.}
Assume $\ell\ge5$ and the original $\mathrm{TopCore}$ for the
embedding list, uniformizer and local inertia subgroup of C80L1.
On the rank block $\mathrm{SetC}$ gives its $\mathrm{TopHyp}$ with
representation $P=\theta M$, lattice $T$ and characters $\alpha_c$.
Use $X_a$ of C80L2 as its finite coefficient module.
Untwisting by $\alpha_1^{-1}=\theta^{-1}$ gives the representation
$M$ on the same vector space and the same underlying integral
lattice, as required by the rank clause of $\mathrm{TopCore}$.
By C80L4 that clause puts $z_a$ in $\mathrm{InHfX}$.
Explicitly there are finitely many integral cocycles $c_j$ which
are rationally Bloch--Kato finite for the untwisted representation,
$x_j\in X_a$, and $m\in T\otimes X_a$ with
\[
 z_a(g)=\sum_j c_j(g)\otimes x_j+
                \theta(g)^{-1}(P(g)\otimes1)m-m .
\]
Choose any lifts $s_j\in L$ of $x_j\in L/O$. Applying $\iota$
gives the reduction of the rational cocycle $\sum_j s_j c_j(g)$
plus the coboundary $(M(g)-1)\iota(m)$. The sum is still
Bloch--Kato finite: linear combinations of the rational crystalline
cobounding primitives give a primitive for the sum after tensoring
with $B_{\rm cris}$. Reduction of such a rational cocycle and addition
of a torsion coboundary is precisely the quotient finite condition
$H_f^1$ in its cocycle formulation. Thus the enlargement of the
original rank-block cut cocycle belongs to the required finite part.
% END C80L5

% BEGIN C80L6
\paragraph{C80L6. The character blocks and the full enlarged cut theorem.}
Keep the source hypotheses: a good continuous family, $\pi\in Z_i$
with the given semisimple splitting, the specified place
$w\mid\ell$, a coefficient enlargement satisfying the given
$\mathrm{CaseA}$, its irreducible $\lambda$, the shift and cyclic
tables satisfying $\mathrm{SetC}$, and $\mathrm{TopCore}$ at
$\ell\ge5$. At every $2\le c\le t$ use C80L2--C80L4 with
the standard integral character line for $\alpha_c$. The second
clause of $\mathrm{TopCore}$ gives a vector $m$ with
$z_a(g)=(\alpha_c(g)/\alpha_1(g)-1)m$. After applying
$\iota$ this is $(\kappa_c(g)-1)\iota(m)$ in the original
enlarged quotient line. Hence the enlarged cut cocycle is an
actual coboundary on the entire local Galois group.

For $c=1$, use that same construction, which also permits the first
character ($t\ge1$ by $\mathrm{SetC}$). The last clause of
$\mathrm{TopCore}$ gives one cobounding vector on inertia for the
ratio $\alpha_1/\alpha_1=1$. Its coboundary is identically zero
there. Applying $\iota$ proves the enlarged first-character
cocycle vanishes on inertia and supplies the zero vector as the
native unramified-cocycle witness. These two assertions and C80L5
prove the full source theorem on the three enlarged cut-cocycle
conditions for every element of each specified cut set. S80J6
provides the stated shift on the chosen family. L79C1--L79C2
then descend them to the original coefficient field at this place.
% END C80L6
